\documentclass[11pt,a4paper]{article}
\usepackage[margin=27mm]{geometry}
\usepackage{amsmath,amssymb,amsthm,mathtools}
\usepackage{graphicx}
\usepackage{xcolor}
\usepackage{microtype}
\usepackage[colorlinks=true,linkcolor=blue!55!black,urlcolor=blue!55!black,citecolor=blue!55!black]{hyperref}

\newtheorem{theorem}{Theorem}
\newcommand{\N}{\mathbb N}
\newcommand{\I}{\mathcal I}

\title{No subset need simultaneously frame\\an unconditional Schauder frame and its coefficient vectors}
\author{Candidate full counterexample to a question in arXiv:2602.21616}
\date{11 August 2026}

\begin{document}
\maketitle

\begin{center}
\fcolorbox{blue!55!black}{blue!5}{\parbox{0.91\textwidth}{
\textbf{Status: candidate full counterexample; likely valid; human review needed.}
There is a normalized unconditional Schauder frame $(x_n,y_n)$, with every
$y_n\ne0$, for which no subset $J$ makes both
$\{x_n/\lVert x_n\rVert\}_{n\in J}$ and
$\{\lVert x_n\rVert y_n\}_{n\in J}$ frames.}}
\end{center}

\section*{Source question}
Yu considers an unconditional Schauder frame $\{x_n\}$ with associated
coefficient vectors $\{y_n\}$, so that
\[
 z=\sum_n\langle z,y_n\rangle x_n
\]
unconditionally for every $z$ in the Hilbert space.  Remark~3.4(b) asks
whether there must be a subset $J$ for which both the normalized primal
sequence and the correspondingly rescaled coefficient sequence are frames
\cite[p.~12]{Yu}.  The exact source text follows.

\begin{center}
\includegraphics[width=0.97\textwidth]{figures/open_problem_crop.png}
\end{center}

\section*{Counterexample and proof intuition}
\begin{theorem}\label{thm:main}
Let $H=\ell_2(\N)$.  There are sequences $(x_n)$ and $(y_n)$ in $H$ such
that:
\begin{enumerate}
 \item $\lVert x_n\rVert=1$ and $y_n\ne0$ for every $n$;
 \item $z=\sum_n\langle z,y_n\rangle x_n$ unconditionally for every $z\in H$;
 \item no subset $J\subseteq\N$ makes both
 \[
   \left\{\frac{x_n}{\lVert x_n\rVert}\right\}_{n\in J}
   \quad\text{and}\quad
   \{\lVert x_n\rVert y_n\}_{n\in J}
 \]
 frames for $H$.
\end{enumerate}
Consequently the question in Remark~3.4(b) has a negative answer.
\end{theorem}

\paragraph{Proof intuition.}
Split the reconstruction of the $k$th coordinate equally among $k$ identical
copies.  The primal sequence then has $k$ copies of the unit vector $e_k$,
while every corresponding coefficient vector has size $1/k$.  If a subset
selects $r_k$ copies from the $k$th block, an upper frame bound for the primal
sequence forces $r_k$ to remain bounded.  A lower frame bound for the
coefficient sequence would instead require $r_k/k^2$ to stay bounded away
from zero.  These two requirements are incompatible as $k\to\infty$.

\section*{Proof}
Let $(e_k)_{k\ge1}$ be the standard orthonormal basis of $H$, and use the
countable index set
\[
 \I=\{(k,j): k\ge1,\ 1\le j\le k\}.
\]
(Fixing any bijection $\I\to\N$ converts the construction into ordinary
sequence notation.)  For $(k,j)\in\I$, set
\begin{equation}\label{eq:def}
 x_{k,j}=e_k,
 \qquad
 y_{k,j}=\frac1k e_k.
\end{equation}
Thus all primal vectors have norm one and all coefficient vectors are
nonzero.

If $z=\sum_{k\ge1}z_ke_k$, then the complete $k$th block contributes
\[
 \sum_{j=1}^k\langle z,y_{k,j}\rangle x_{k,j}
 =\sum_{j=1}^k\frac{z_k}{k}e_k=z_ke_k.
\]
This reconstruction is unconditional, not merely conditionally convergent.
Indeed, for a finite $F\subset\I$ put
\[
 r_k(F)=\#\{j:(k,j)\in F\},
 \qquad 0\le r_k(F)\le k.
\]
The corresponding partial sum is the diagonal contraction
\begin{equation}\label{eq:partial}
 S_Fz=\sum_{(k,j)\in F}\langle z,y_{k,j}\rangle x_{k,j}
     =\sum_{k\ge1}\frac{r_k(F)}k z_ke_k.
\end{equation}
Given $\varepsilon>0$, choose $K$ so that
$\sum_{k>K}|z_k|^2<\varepsilon^2$, and let $F_0$ contain every index in the
first $K$ blocks.  For every finite $F\supseteq F_0$, equation
\eqref{eq:partial} gives
\[
 \lVert z-S_Fz\rVert^2
 =\sum_{k>K}\left|1-\frac{r_k(F)}k\right|^2|z_k|^2
 \le \sum_{k>K}|z_k|^2<\varepsilon^2.
\]
Hence the net of finite partial sums converges to $z$, which is the
unconditional convergence criterion.  Therefore \eqref{eq:def} is a
normalized unconditional Schauder frame with the displayed associated
coefficient vectors.

Now let $J\subseteq\I$ be arbitrary and define its selected multiplicities
\[
 r_k=\#\{j:(k,j)\in J\},\qquad 0\le r_k\le k.
\]
Since $\lVert x_{k,j}\rVert=1$, the two frame quadratic forms in the source
question are, exactly,
\begin{align}
 \sum_{(k,j)\in J}
  \left|\left\langle z,\frac{x_{k,j}}{\lVert x_{k,j}\rVert}\right\rangle\right|^2
 &=\sum_{k\ge1}r_k|z_k|^2,\label{eq:primal}\\
 \sum_{(k,j)\in J}
  |\langle z,\lVert x_{k,j}\rVert y_{k,j}\rangle|^2
 &=\sum_{k\ge1}\frac{r_k}{k^2}|z_k|^2.\label{eq:dual}
\end{align}
Suppose first that the sequence in \eqref{eq:primal} is a frame, and let
$B<\infty$ be an upper frame bound.  Testing on $z=e_k$ yields
\begin{equation}\label{eq:boundedmultiplicity}
 r_k\le B\qquad(k\ge1).
\end{equation}
If the sequence in \eqref{eq:dual} also had a positive lower frame bound
$A'$, testing again on $e_k$ would give
\[
 A'\le \frac{r_k}{k^2}\le\frac{B}{k^2}\qquad(k\ge1),
\]
which is impossible as $k\to\infty$.  Thus no $J$ has the two required frame
properties, proving the theorem. \qed

\section*{Scope and novelty note}
The counterexample answers only the simultaneous-subset question in
Remark~3.4(b); it does not contradict the paper's theorem that the normalized
primal sequence contains a frame subsequence.  In fact, choosing one copy of
each $e_k$ gives such a primal orthonormal basis, while its paired coefficient
vectors are $e_k/k$ and have no lower frame bound.

The run's four cheap indexes and bounded searches through 11 August 2026 for
the exact remark and its central terms found the source paper but no later
resolution.  Because the source is recent, novelty remains provisional
pending expert literature review.

\begin{thebibliography}{9}
\bibitem{Yu}
P.-T.~Yu,
\emph{Every semi-normalized unconditional Schauder frame in Hilbert spaces
contains a frame}, arXiv:2602.21616v3 (2026).
\end{thebibliography}

\end{document}
